\subsection{从李代数到李群的过渡}

定理 3. (i) 若 L 是有限维李代数，则存在一个单连通李群 G，使 L(G) 同构于 L。

\begin{enumerate}
\def\labelenumi{(\roman{enumi})}
\setcounter{enumi}{1}
\tightlist
\item
  设 \(G_{1}\) 和 \(G_{2}\) 是两个连通李群，且 \(G_{1}\) 单连通。令 f 为从 \(\mathrm{L}(\mathrm{G}_{1})\) 到 \(\mathrm{L}(\mathrm{G}_{2})\) 的同构，\(\phi\) 为从 \(G_{1}\) 到 \(G_{2}\) 的态射，满足 \(\mathrm{L}(\phi)=f\)，N 为 \(\phi\) 的核。则 N 是 \(G_{1}\) 的中心中的离散子群，而从 \(G_{1}/N\) 到 \(G_{2}\) 的态射由 \(\phi\) 通过商得到，且是李群同构。若 \(G_{2}\) 单连通，则 \(\phi\) 是同构。
\end{enumerate}

令 \(\mathbf{L}\) 为有限维李代数。存在有限维向量空间 \(\mathbf{E}\)，使得 \(\mathbf{L}\) 可以识别为 \(\operatorname{End}(\mathbf{E})\) 的李子代数（Chapter I, §7, Theorem 2）。令 \(\mathbf{H}\) 为 \(\mathbf{GL}(\mathbf{E})\) 中李代数为

L 的积分子群。令 \(\hat{\mathbf{H}}\) 为它的万有覆盖（§ 1, no. 9, Remark）。则 \(\mathbf{L}(\hat{\mathbf{H}})\) 同构于 L，从而得到 (i)。

  令 \(G_{1}, G_{2}, \tilde{f}, \phi, N\) 如 (ii) 中所述。则 \(\phi\) 是 étale 的，因而 \(\phi(G_{1})\) 是 \(G_{2}\) 的开子群；于是 \(\phi(G_{1}) = G_{2}\)。另一方面，N 是离散的，因而包含于 \(G_{1}\) 的中心中（Integration, Chapter VII, § 3, Lemma 4）。显然，从 \(G_{1}/N\) 到 \(G_{2}\) 的态射由 \(\phi\) 通过商得到，且是李群同构。若 \(G_{2}\) 单连通，则每个从 \(G_{1}\) 到 \(G_{2}\) 的 étale 映射都是单射，因而 \(N = \{e\}\)。

命题 5. 设 G 是连通实李群。假设 L(G) 带有一个复可赋范李代数结构 L'，与其可赋范实李代数结构相容。G 上存在唯一一个复李群结构，该结构与实李群结构相容，且李代数为 L'。

由 § 4, no. 2, Corollary 2 to Theorem 2，只需证明 \(L'\) 的结构在 Ad G 下不变。令 \(\phi\) 为 G 的一个指数映射。由 § 4, no. 4, Corollary 3 (i) to Proposition 8，存在 \(L(G)\) 中 0 的一个邻域 V，使 \(L'\) 的结构在 \(Ad\phi(V)\) 下不变。但是 \(Ad\phi(V)\) 生成 G，因为 G 是连通的。

若不假设 G 连通，命题 5 的结论不一定成立（练习 7）。

命题 6. 设 G 是连通复李群。若 G 紧，则 G 交换。

全纯映射 \(g \mapsto \operatorname{Ad} g\) 从 \(\mathbf{G}\) 到 \(\mathcal{L}(\mathbf{L}(\mathbf{G}))\) 是常值的（Differentiable and Analytic Manifolds, R, 3.3.7），因而 ad \(a = 0\) 对所有 \(a \in \mathbf{L}(\mathbf{G})\) 成立（§ 3, no. 12, Proposition 44）。所以 \(\mathbf{G}\) 是交换的（§ 4, Corollary 3 to Theorem 1）。

